\chapter{Grundregeln}

Dieses Kapitel erklärt alle grundlegenden Spielmechaniken, die du als Spieler
kennen musst. Wenn im Spiel etwas unklar oder riskant ist, greifen diese Regeln.

\section{Proben}

\begin{rulebox}{Was ist eine Probe?}
	Eine \term{Probe}{probe} wird immer dann gewürfelt, wenn der Ausgang einer
	Handlung ungewiss ist und ihr Erfolg oder Misserfolg Konsequenzen hat.
\end{rulebox}

Eine Probe besteht aus:
\begin{itemize}
	\item einem \textbf{Würfel} (abhängig vom RANG)
	\item möglichen \textbf{Modifikatoren}
	\item deiner \term{Wurfzahl}{wurfzahl} (Wurf + Modifikatoren)
	\item einem vom Spielleiter festgelegten \term{Zielwert}{zielwert}
\end{itemize}

\begin{rulebox}{Wurfzahl und Zielwert}
	Die \textbf{Wurfzahl} ist dein Endergebnis nach dem Würfeln inklusive aller
	Modifikatoren.  
	
	Der \textbf{Zielwert} beschreibt die Schwierigkeit der Handlung und wird vom
	Spielleiter festgelegt.
\end{rulebox}

Eine Probe ist erfolgreich, wenn:
\[
\text{Wurfzahl} \ge \text{Zielwert}
\]

\begin{rulebox}{Standard-Zielwert}
	Wenn keine besondere Schwierigkeit vorliegt, gilt als Standard:
	\textbf{Zielwert 4}.
\end{rulebox}

\begin{examplebox}{Einfache Probe}
	Du willst über eine niedrige Mauer klettern.  
	Der Spielleiter legt einen Zielwert von 6 fest.  
	
	Du würfelst eine 7 (Wurfzahl 7) – die Aktion gelingt.
\end{examplebox}

\section{Rang und Würfel}

\begin{rulebox}{Was bedeutet Rang?}
	Der \term{Rang}{rang} beschreibt, wie gut dein Charakter in einer Fähigkeit,
	einem Skill oder einem Zauber ist.
\end{rulebox}

Jeder Rang ist fest einem Würfeltyp zugeordnet:

\begin{center}
	\begin{tabular}{lll}
		\toprule
		\textbf{RANG} & \textbf{Würfel} & \textbf{Chance auf Wurfzahl $\ge 4$} \\
		\midrule
		1 & W4  & 25\% \\
		2 & W6  & 50\% \\
		3 & W8  & 62{,}5\% \\
		4 & W10 & 70\% \\
		5 & W12 & 75\% \\
		6 & W20 & 85\% \\
		\bottomrule
	\end{tabular}
\end{center}

Je höher dein Rang, desto höher die maximale Augenzahl und damit die
Erfolgschance.

\section{Die Ass-Regel (Explodierende Würfel)}

\begin{rulebox}{\term{Ass-Regel}{ass-regel}}
	Wenn du bei einem Wurf die \textbf{höchste mögliche Zahl} deines Würfels
	erreichst, wird der Würfel erneut geworfen und das Ergebnis addiert.
\end{rulebox}

Dieser Vorgang kann sich mehrfach wiederholen.

\begin{examplebox}{Ass-Regel in Aktion}
	Du würfelst einen W8 und erzielst eine 8.  
	Du würfelst erneut und bekommst eine 5.  
	Das Gesamtergebnis ist \textbf{13}.
\end{examplebox}

Die Ass-Regel sorgt dafür, dass auch sehr schwierige Proben theoretisch
möglich bleiben.

\section{Modifikatoren: Boni und Mali}

\begin{rulebox}{Modifikatoren}
	Ein \term{Bonus}{bonus} oder \term{Malus}{malus} verändert das Ergebnis eines Wurfs positiv oder negativ.
\end{rulebox}

Typische Quellen für Modifikatoren sind:
\begin{itemize}
	\item situative Entscheidungen des Spielleiters (z.\,B. Vorteil/Nachteil durch die Szene)
	\item Ausrüstung
	\item Umstände (Zeitdruck, Dunkelheit, Verletzungen)
	\item aktive Effekte durch Fähigkeiten
\end{itemize}

Modifikatoren werden immer \textbf{nach} dem Wurf angewendet.


\section{Kombinationswürfe}

\begin{rulebox}{\term{Kombinationswürfe}{kombinationswurf}}
	Manche Handlungen erlauben es, mehrere passende Talente zu kombinieren.
\end{rulebox}

Dabei gilt:
\begin{itemize}
	\item Für die Probe wird der \textbf{höchste beteiligte Rang} verwendet.
	\item Jedes \textbf{weitere passende Talent} gibt einen Bonus auf die Wurfzahl:
	\begin{itemize}
		\item \textbf{+1}, wenn dieses Talent mindestens Rang 3 hat
		\item \textbf{+2}, wenn dieses Talent mindestens Rang 5 hat
	\end{itemize}
\end{itemize}

Die genaue Zulässigkeit entscheidet der Spielleiter.

\begin{examplebox}{Kombinationswurf}
	Du versuchst, ein Gespräch über Kriegsstrategie zu führen und kombinierst
	\textbf{Kampf Rang 4} mit \textbf{Wissen Rang 3}.
	Du würfelst mit Rang 4 und erhältst zusätzlich \textbf{+1} durch Wissen.
\end{examplebox}


\section{Erfolge und Fehlschläge}

\begin{rulebox}{Erfolg und Fehlschlag}
	Der Ausgang einer Probe ergibt sich aus dem Vergleich von
	Wurfzahl und Zielwert.
\end{rulebox}

\textbf{Grundregel:}
\begin{itemize}
	\item Wird der \textbf{Zielwert nicht erreicht}, ist die Probe ein \textbf{Fehlschlag}.
	\item Wird der \textbf{Zielwert erreicht oder übertroffen}, ist die Probe ein \textbf{Erfolg}.
\end{itemize}

Der \textbf{Standard-Zielwert} für Proben beträgt \textbf{4}, sofern keine besondere
Schwierigkeit vorliegt.

\medskip

\textbf{Qualität des Ergebnisses:}  
Je höher die \textbf{Wurfzahl} (Wurf + Modifikatoren), desto besser fällt das
Ergebnis der Handlung aus.

\begin{rulebox}{Konsequenzen}
	Es gibt keine festen Kategorien für kritische Erfolge oder Fehlschläge.
	Die Konsequenzen einer Probe ergeben sich immer aus dem \textbf{Kontext der Szene}.
\end{rulebox}

Ein Fehlschlag kann je nach Situation unterschiedliche Folgen haben:
\begin{itemize}
	\item Bei einer Heimlichkeitsprobe kann ein Fehlschlag Alarm auslösen.
	\item Bei einem Armdrücken kann ein Fehlschlag lediglich den Verlust von Gold bedeuten.
\end{itemize}

Der Spielleiter entscheidet situativ über Ausmaß und Art der Konsequenzen.


\section{Wann wird nicht gewürfelt?}

\begin{rulebox}{Kein Wurf nötig}
	Wenn eine Handlung trivial ist oder keine Konsequenzen hat, wird \textbf{nicht}
	gewürfelt.
\end{rulebox}

Ein erfahrener Krieger fällt nicht vom Pferd, nur weil er reitet.  
Ein Magier kennt einfache arkane Grundregeln.

Würfe sind nur dann relevant, wenn sie das Spiel voranbringen.

